\section{The near-miss mechanism step by step (full derivation of Section~\ref{sec:string})}\label{app:derivation}

Section~\ref{sec:string} gives the results of the P1 theory in condensed form. This appendix gives every step, in the
order in which the argument is built, so that it can be checked line by line. It follows the Stage-1 note
\rep{analytic/NEAR_MISS_IMPULSE.md} (v5.1); each step names the sympy or numerical check that verifies it
(\rep{analytic/checks/cNN_*.py}, with outputs next to them). Units are $\hbar=c=1$; the notation is that of
Appendix~\ref{app:notation}. Table~\ref{tab:derivation_status} at the end lists which steps are derived, which are
verified numerically and which are numerical only.

\subsection{Equations of motion}\label{app:der:eom}
The system is a classical Nambu--Goto (NG) string of tension $\sigma$ whose two ends carry point masses $\mQ$~\cite{Kang:2008ea}. Here the ends also carry
charges $q_1=+q$, $q_2=-q$ and sit in a uniform constant external field:
\begin{equation}
  S=-\sigma\int d^2\xi\,\sqrt{-\det h_{ab}}\;-\;\sum_{i=1,2}\mQ\int d\tau_i\;+\;\sum_{i=1,2}q_i\int A_\mu\,dx_i^\mu ,
  \label{eq:D:action}
\end{equation}
where $h_{ab}$ is the induced metric on the world sheet and $\tau_i$ the proper time of end $i$. Two consequences are
used throughout.
\begin{enumerate}
\item \textbf{The string is neutral.} The field acts only on the ends. The string exchanges momentum and angular momentum
  with the ends but conserves the totals, so the total momentum and the total angular momentum (ends plus string) change
  only through the Lorentz forces on the ends.
\item \textbf{Bulk and boundary.} For a string lying along $x$ with a small transverse displacement $y(t,x)$, the static
  gauge $X^\mu=(t,x,y(t,x))$ gives the Lagrangian density $-\sigma\sqrt{1+y_x^2-y_t^2}\simeq-\sigma+\tfrac{\sigma}{2}(y_t^2-y_x^2)$.
  So $y$ obeys the wave equation $y_{tt}=y_{xx}$ and the string carries the transverse momentum density $\sigma y_t$.
  The force that the string exerts on an end follows from momentum balance (Section~\ref{app:der:wedge}). For a straight
  string whose end moves along it, the force is $-\sigma\hat s$, of magnitude $\sigma$ and directed toward the other end.
\end{enumerate}
The infracolour gauge group enters only through $\sigma=\Lam^2$. All results below therefore hold for any $N_{\rm IC}$
and representation once $\sigma$ is fixed.

\subsection{From the lab solenoid to the pair rest frame}\label{app:der:frame}
In the lab there is only a magnetic field $\mathbf B=B\hat z$. Let $\boldsymbol\beta_{\rm cm}$ be the initial velocity of
the pair and $\gamma_{\rm cm}$ its Lorentz factor. The standard field transformation (e.g.\ Jackson, Sec.~11.10) gives,
in the frame moving with $\boldsymbol\beta_{\rm cm}$ (``CM frame''),
\begin{equation}
  \Estar=\gamma_{\rm cm}\,\boldsymbol\beta_{\rm cm}\times\mathbf B,\qquad
  \Bstar=\gamma_{\rm cm}\mathbf B-(\gamma_{\rm cm}-1)(\mathbf B\cdot\hat\beta_{\rm cm})\,\hat\beta_{\rm cm}.
  \label{eq:D:fields}
\end{equation}
Both fields are uniform and constant, and $|\Estar|=\gamma_{\rm cm}\beta_T B$, with $\beta_T$ the component of
$\boldsymbol\beta_{\rm cm}$ transverse to $\hat z$. The action~\eqref{eq:D:action} is Lorentz invariant. One can therefore
solve the problem in the CM frame and boost the hits back to the lab, and no simultaneity convention enters. In the CM
frame the two ends start at the origin with momenta $\pm p^*\hat n$. Here $\hat n$ is the direction of end~1 (charge $+q$)
and $p^*=\mQ\rho$ with $\rho\equiv\gst\bst$. The pair invariant mass is $M=2\mQ\gst$, so
\begin{equation}
  \gamma_{\rm cm}\beta_T=\frac{p_T}{M}=\frac{p_T}{2\mQ\gst}.
  \label{eq:D:gbT}
\end{equation}
Split $\Estar$ into its components along $\hat n$ and perpendicular to it, $\Estar_\perp=|\Estar_\perp|\hat e_\perp$,
and let $\alpha$ be the angle between $\Estar$ and $\hat n$. The dimensionless field strengths of Eq.~\eqref{eq:eps} are
then
\begin{equation}
  \epsB=\frac{qB}{\sigma},\qquad
  \epsp=\frac{q|\Estar_\perp|}{\sigma}=\epsB\,\gamma_{\rm cm}\beta_T\sin\alpha,\qquad
  \kB=\frac{|\Bstar_\perp|}{|\Estar_\perp|}.
  \label{eq:D:eps}
\end{equation}
With $qB=2.998\times10^{-3}\GeV/\mathrm{cm}$ per tesla and $\sigma=50.68\GeV/\mathrm{cm}\,(\Lam/\keV)^2$, one gets
$\epsB=2.254\times10^{-4}$ at $3.8112\tesla$ and $1\keV$. Typical events have $\epsp\sim10^{-5}$--$10^{-4}$ at $1\keV$,
so the field is a very weak perturbation of the string force.

\subsection{The vacuum yo-yo (zeroth order)}\label{app:der:yoyo}
\paragraph{Solution.} Without fields the string stays straight along $\hat n$ and each end moves along it. The force on
each end is therefore exactly $\sigma$, directed toward the other end~\cite[Sec.~3]{Kang:2008ea}. For end~1,
$dp_1/dt=-\sigma$, so $p_1(t)=p^*-\sigma t$ and
\begin{equation}
  x_1(t)=\int_0^t\frac{p_1\,dt'}{\sqrt{\mQ^2+p_1^2}}=\frac{1}{\sigma}\Bigl[\sqrt{\mQ^2+p^{*2}}-\sqrt{\mQ^2+(p^*-\sigma t)^2}\Bigr].
\end{equation}
The separation $x=x_1-x_2=2x_1$ is therefore~\cite{Kang:2008ea} (check \rep{c01} (1))
\begin{equation}
  x(t)=2\lzero\Bigl[\gst-\sqrt{1+(\rho-t/\lzero)^2}\Bigr],\qquad 0\le t\le T_h=\frac{2p^*}{\sigma}=2\rho\lzero,
  \qquad \lzero=\frac{\mQ}{\sigma}.
  \label{eq:D:x}
\end{equation}
It vanishes at $t=0$ and $t=T_h$, reaches $\lmax=2(\gst-1)\lzero$ at $t=\rho\lzero$, and the ends pass each other at
speed $\bst$. At $t=T_h$ the ends pass \emph{through} each other, and the next half-period is the mirror image,
$x(t+T_h)=-x(t)$. The full CM period is $2T_h=4p^*/\sigma$. We number the crossings $j=1,2,\ldots$: odd $j$ at
$t=T_h,3T_h,\ldots$, even $j$ at $t=2T_h,4T_h,\ldots$ (production is the crossing $j=0$).

\paragraph{Zero-length string at the crossings.} The world sheet is a chain of flat ``lenses'' bounded by the two end
world lines, which touch at the crossings. At a crossing the string has zero length, and hence zero energy, momentum and
angular momentum. The whole state is given by the two end momenta at a common point.

\paragraph{Half-period integral.} The torque in Section~\ref{app:der:torque} needs $I\equiv\int_0^{T_h}x\,dt$.
Substitute $u=\rho-t/\lzero$ and use $\sqrt{1+\rho^2}=\gst$ (check \rep{c01} (2)):
\begin{equation}
  I=2\lzero^2\int_{-\rho}^{\rho}\bigl[\gst-\sqrt{1+u^2}\bigr]du
   =2\lzero^2\Bigl[2\rho\gst-\bigl(\rho\sqrt{1+\rho^2}+\operatorname{asinh}\rho\bigr)\Bigr]
   =2\lzero^2\bigl[\gst\rho-\operatorname{asinh}\rho\bigr].
  \label{eq:D:I}
\end{equation}

\paragraph{Rigid-rotation lemma.} \emph{A transverse momentum $\Delta p_\perp$ given to the relative motion around a crossing
rotates the whole later motion rigidly by $\delta=\Delta p_\perp/p^*$.} Proof: at a crossing the string has zero length
(previous paragraph), so the state is two point momenta at one point. Without fields the equations are rotation
invariant, so the rotated initial state evolves into the rotated vacuum solution. A kick changes the momenta from
$\pm p^*\hat n$ to $\pm(p^*\hat n+\Delta p_\perp\hat e)$, which equals $\pm p^*(\hat n+\delta\hat e)$ up to a change of
magnitude of $O(\delta^2)$. The subsequent motion is the yo-yo along the axis $\hat n+\delta\hat e$. In other words, the
axis direction is a zero mode of the yo-yo. Anything that happens in a short window around a crossing acts on the
long-time motion only through this zero mode.

\subsection{Torque, angular momentum and the miss distance}\label{app:der:torque}
\paragraph{Torque law.} Let $\mathbf J=\sum_i\mathbf x_i\times\mathbf p_i+\mathbf J_{\rm string}$ be the total angular
momentum about the CM. By item~1 of Section~\ref{app:der:eom} it changes only through the Lorentz forces on the ends:
\begin{equation}
  \frac{d\mathbf J}{dt}=\sum_i\mathbf x_i\times q_i\bigl(\Estar+\mathbf v_i\times\Bstar\bigr).
\end{equation}
To first order in the field, evaluate the right-hand side on the yo-yo, where $\mathbf x_2=-\mathbf x_1$,
$\mathbf v_2=-\mathbf v_1$ and $q_2=-q_1$. The magnetic terms cancel exactly:
$\mathbf x_1\times q(\mathbf v_1\times\Bstar)+(-\mathbf x_1)\times(-q)\bigl((-\mathbf v_1)\times\Bstar\bigr)=0$. Their
contribution to the relative force, $(\mathbf F_1-\mathbf F_2)/2$, also vanishes (check \rep{c01} (4)). The electric term
gives
\begin{equation}
  \frac{d\mathbf J}{dt}=q(\mathbf x_1-\mathbf x_2)\times\Estar=q\,x(t)\,\hat n\times\Estar
  =q|\Estar_\perp|\,x(t)\,\hat n\times\hat e_\perp .
  \label{eq:D:torque}
\end{equation}
Only $\Estar_\perp$ produces a torque, and in linear response $\Bstar$ does nothing to the relative motion. The
symmetry reason is that under inversion through the CM combined with exchange of the ends, the $\Estar$ force is even
and the $\Bstar$ force is odd.

\paragraph{The alternating pattern.} Integrating Eq.~\eqref{eq:D:torque} over the first half-period gives
$\mathbf J(T_h)=q|\Estar_\perp|\,I\,\hat n\times\hat e_\perp$. Over the second half-period $x$ changes sign, so
$\mathbf J(2T_h)=0$. The angular momentum is therefore zero at production and at every even crossing, and equal to
$\mathbf J_{\rm odd}=q|\Estar_\perp|I\,\hat n\times\hat e_\perp$ at every odd crossing (Li et al.'s ``meet at
$g=2n$''~\cite[Eqs.~(A.4)--(A.5)]{Li:2020aoq}, here with the value at the odd crossings).

\paragraph{Miss distance.} At an odd crossing the string has zero length, so $\mathbf J$ is carried by the ends. If they
pass at transverse offsets $\pm(b/2)\hat s$ with momenta $\mp p^*\hat n$ (end~1 is now moving back),
$\mathbf J=\tfrac b2\hat s\times(-p^*\hat n)+(-\tfrac b2\hat s)\times(p^*\hat n)=b\,p^*\,\hat n\times\hat s$. Comparing
with $\mathbf J_{\rm odd}$ gives $\hat s=\hat e_\perp$: the positive end passes on the $+\Estar_\perp$ side, as
expected. The miss distance is $b=J/p^*$. With Eq.~\eqref{eq:D:I}, $p^*=\mQ\rho$ and $\lzero^2/\mQ=\lzero/\sigma$
(check \rep{c01} (3)):
\begin{equation}
  \bodd=\frac{q|\Estar_\perp|\,2\lzero^2(\gst\rho-\operatorname{asinh}\rho)}{\mQ\rho}
  =2\epsp\lzero\Bigl[\gst-\frac{\operatorname{asinh}\rho}{\rho}\Bigr],\qquad
  \frac{\bodd}{\epsp\lmax}=\frac{\gst-\operatorname{asinh}\rho/\rho}{\gst-1}.
  \label{eq:D:bodd}
\end{equation}
This is Eq.~\eqref{eq:bodd}. For slow ends, $\gst\simeq1+\rho^2/2$ and $\operatorname{asinh}\rho/\rho\simeq1-\rho^2/6$,
so $\gst-\operatorname{asinh}\rho/\rho\simeq\tfrac23\rho^2$ and $\bodd\to\tfrac43\epsp\lzero{\bst}^2$, i.e.\
$\bodd/(\epsp\lmax)\to4/3$. The ratio is 1.327/1.298/1.240/1.188 at $\gst=1.05/1.3/2/3$. The causal solver reproduces it to $10^{-4}$
(check \rep{c03}); two solvers (causal and conformal) agree with it to $\le10^{-3}$ (Table~\ref{tab:numerics}). Even crossings have $b=0$ at this order. At $500\GeV$, $1\keV$ a
typical event has $\epsp\approx5\times10^{-5}$ and $\lmax\approx4.5\cm$, so $\bodd\approx3\um$, against a median of
$\approx3.5\um\,(\keV/\Lam)^4$ over the sample (the $\Lam^{-4}$ comes from $\epsp\propto\Lam^{-2}$ and $\lzero\propto\Lam^{-2}$).

\subsection{Linear response and where it fails}\label{app:der:linear}
The failure is easiest to see for a straight string with slow ends. Let $\mathbf r=\mathbf x_1-\mathbf x_2=x\hat
n+z\hat e_\perp$, with reduced mass $\mu=\mQ/2$. Each end is pulled toward the other with force $\sigma$, so
$\mu\ddot{\mathbf r}=-\sigma\hat{\mathbf r}+q\Estar$. Its transverse component, to first order in $z$, is
\begin{equation}
  \mu\ddot z=-\sigma\frac{z}{|x(t)|}+q|\Estar_\perp| .
  \label{eq:D:zlin}
\end{equation}
Three exact properties follow (the third is used in check \rep{c05}).
\begin{enumerate}
\item $z=\theta_0x(t)$ solves the homogeneous equation, because $\mu\ddot x=-\sigma\,{\rm sgn}\,x$. This is the
  rotation zero mode.
\item The Wronskian $W=\mu(x\dot z-z\dot x)$, which is the angular momentum $J$, obeys
  $\dot W=\mu(x\ddot z-z\ddot x)=q|\Estar_\perp|\,x$, i.e.\ the torque law~\eqref{eq:D:torque}.
\item The angle $\theta=z/x$ of the relative vector to the original axis obeys $\dot\theta=J/(\mu x^2)$.
\end{enumerate}
Near a crossing at time $s=0$ the separation is $|x|\simeq2v|s|$ (each end moves with $v=\bst$), so the coefficient
$\sigma/|x|\propto1/|s|$ makes $s=0$ a regular singular point of Eq.~\eqref{eq:D:zlin}. Let $z_c$ be the offset at the
crossing. Integrating through it, $\mu\Delta\dot z\simeq-\int\sigma z_c\,ds/(2v|s|)$ diverges logarithmically unless
$z_c=0$. Thus $z$ stays continuous but $\dot z$ has a $\ln|s|$ singularity.
\begin{itemize}
\item At production and at even crossings, $z_c=J/p^*=0$. The linear solution is regular, and first-order perturbation
  theory is exact there. Numerically, the response before the first crossing is linear in $B$ to $10^{-4}$.
\item At odd crossings, $z_c=\bodd\ne0$. The linear equation then has no unique continuation through the crossing. The
  outcome is set by a region $|x|\lesssim\bodd$ in which the small-angle assumption $|z|\ll|x|$ fails, and the impulse
  is $\propto\bodd\ln(1/\bodd)$, which is not linear in the field.
\end{itemize}
So first-order perturbation theory fails at every odd crossing and only there. Physically, the odd crossing is a
particle passing a constant-magnitude central force at impact parameter $\bodd$. The logarithm of the deflection is
elementary orbit mechanics (Section~\ref{app:der:nr}); what is specific to quirks is where it occurs and how large $b$ is.
Li et al.\ noted that the first-order solution fails on the approach to the crossing~\cite[App.~A]{Li:2020aoq}.

\subsection{The intermediate region: an exact NG solution with retardation}\label{app:der:wedge}
Use the time $s$ from the odd crossing and look at the region $\bodd\ll|x|\ll\lmax$, i.e.\ $|s|\ll T_h$. There:
\begin{itemize}
\item the ends move with nearly constant velocities $\pm v$, $v=\bst$, at $x=\pm vs$ (the deceleration changes $v$ by
  $O(\sigma|s|/p^*)$);
\item the whole string between them has length $2v|s|$ and energy $2\sigma v|s|\ll\mQ\gst$, so it is light and slaved to
  the motion of its ends;
\item in static gauge the unperturbed world sheet is the flat wedge $|x|<v|s|$, and the transverse displacement obeys
  $y_{ss}=y_{xx}$ there (Section~\ref{app:der:eom});
\item the ends are offset by $\pm b/2$, with $b=\bodd$: $y(s,\pm vs)=\pm b/2$.
\end{itemize}

\paragraph{Solution.} Apart from $b$, which only sets the amplitude, the problem has no scale. Try a self-similar
solution $y=f(u)$ with $u=x/s$. Then $y_{xx}=f''/s^2$ and $y_{ss}=(u^2f''+2uf')/s^2$, and the wave equation becomes
$(1-u^2)f''=2uf'$. Hence $f'\propto1/(1-u^2)$ and $f\propto\operatorname{artanh}u$. The boundary values fix the
constant (check \rep{c02} (2)):
\begin{equation}
  y(s,x)=\frac b2\,\frac{\operatorname{artanh}(x/s)}{\operatorname{artanh}v}.
  \label{eq:D:wedge}
\end{equation}
The string is curved: it is not the straight line between the ends.

\paragraph{Force on the ends.} The transverse momentum of the string between $X_L(s)$ and $X_R(s)$ is
$P=\int_{X_L}^{X_R}\sigma y_s\,dx$. By Leibniz' rule and $y_{ss}=y_{xx}$,
\begin{equation}
  \frac{dP}{ds}=\sigma\bigl[y_x+V y_s\bigr]_{X_L}^{X_R},\qquad V=\dot X .
\end{equation}
Apply this balance to a short piece of string adjacent to the right end: the momentum that flows through the moving
boundary $X_R$ is taken up by that end, since the total transverse momentum of end plus string changes only through the
external force (Section~\ref{app:der:eom}). So the right end receives
$dp_R/ds=-\sigma(y_x+Vy_s)|_{X_R}$, and the left end the opposite with $X_L$ (check \rep{c02} (1)). At $x=vs$,
Eq.~\eqref{eq:D:wedge} gives $y_x=C/[s(1-v^2)]$ and $y_s=-Cv/[s(1-v^2)]$ with $C=b/(2\operatorname{artanh}v)$, so
$y_x+vy_s=C/s$ and (check \rep{c02} (3))
\begin{equation}
  F_\perp=\mp\frac{\sigma b}{2|s|\operatorname{artanh}v}.
  \label{eq:D:Fperp}
\end{equation}
Each end is pulled toward the side of the other end, with the same sign before and after the crossing. Three remarks:
\begin{itemize}
\item For $v\to0$, Eq.~\eqref{eq:D:Fperp} becomes $\sigma b/(2v|s|)$: the tension times the slope $b/(2v|s|)$ of a
  straight string between the ends, i.e.\ the quasi-static value (check \rep{c02} (4)). The factor
  $v/\operatorname{artanh}v=0.97/0.85/0.66/0.53$ at $\gst=1.05/1.28/2/3$ is the retardation effect: the string cannot
  follow the ends instantaneously. The Kang--Luty straight-string force gives instead $\sigma b(1-v^2)/(2v|s|)$ (check
  \rep{c02} (5); Table in \rep{NEAR_MISS} \S4.4).
\item A relative transverse velocity of the ends gives an $s$-independent force, and the free log-periodic string modes
  give convergent impulses (check \rep{c02} (6)). Only the offset $b$ produces a $1/|s|$ force.
\item The approximations fail at $|x|\gtrsim\lmax$, where the ends decelerate, and at $|x|\lesssim\bodd{\gst}^2/\operatorname{artanh}\bst$,
  where the offset is no longer small compared with the separation. These two scales cut off the logarithm.
\end{itemize}

\paragraph{Impulse and rotation.} Integrated over the intermediate region on both sides of the crossing, each end
receives $\int\sigma b\,ds/(2|s|\operatorname{artanh}v)=(\sigma b/\operatorname{artanh}v)\ln(s_{\rm out}/s_{\rm in})$
(check \rep{c02} (7)). The relative transverse momentum $p_{{\rm rel},\perp}=(p_{1\perp}-p_{2\perp})/2$ changes by
\begin{equation}
  \Delta p_{{\rm rel},\perp}=-\frac{\sigma\bodd}{\operatorname{artanh}\bst}\Bigl[\ln\frac{\lmax}{\bodd}+O(1)\Bigr]\hat e_\perp .
  \label{eq:D:dprel}
\end{equation}
This momentum is acquired logarithmically, mostly at $|x|\ll\lmax$: about 85\% of the logarithm is collected within
$|s|\lesssim0.05\,T_h$, where $x\lesssim0.2\,\lmax$ and the string is short and light. By the rigid-rotation lemma it
therefore rotates the axis by $\delta=|\Delta p_{{\rm rel},\perp}|/p^*$; the remainder from the finite string length is
$O(1)$ and is part of $K$, which the matched calculation of Section~\ref{app:der:nr} treats exactly. The sign works as
follows. End~1 passes on the $+\hat e_\perp$ side and is pulled toward $-\hat e_\perp$ while it moves along $-\hat n$, so
after the crossing it heads along $-(\hat n+\delta\hat e_\perp)$. The axis, as a line, turns from $\hat n$ toward
$+\hat e_\perp$, i.e.\ toward $+\Estar_\perp$. Inserting $\bodd$ from Eq.~\eqref{eq:D:bodd} and $p^*=\mQ\rho=\sigma\lzero\rho$:
\begin{equation}
  \frac{\sigma\bodd}{p^*\operatorname{artanh}\bst}
  =\frac{2\epsp\bigl[\gst-\operatorname{asinh}\rho/\rho\bigr]}{\rho\operatorname{artanh}\bst}
  \equiv\epsp\,G(\gst),
  \label{eq:D:G}
\end{equation}
which is the coefficient of the logarithm in Eq.~\eqref{eq:rot}. For slow ends, $\gst-\operatorname{asinh}\rho/\rho\simeq\tfrac23\rho^2$ and $\operatorname{artanh}\bst\simeq\rho$, so $G\to\tfrac43$. In
exact-NG numerics the slope of $\rot/\epsp$ against $\ln(1/\epsp)$ equals $G$ to $\le0.03\%$ (check \rep{c03}) and
to 0.1--0.2\% in the two-solver campaign (\rep{NEAR_MISS} \S9.1). A quasi-static string ($1/\bst$ instead of
$1/\operatorname{artanh}\bst$) would give a slope 1.03--1.87 times larger.

\paragraph{No logarithm in the momentum at the turning points.} The logarithm ends up in the \emph{direction} of the
axis, not in a lasting transverse momentum. At the next turning point only $J$ survives, and the transverse momentum
measured from turning point to turning point has no logarithm: $\Delta\Pi/(\sigma b/\bst)=\rho\bst/[(\gst-1)(1+(\gst-1)/3)]$,
measured 1.92/1.61/1.15/0.79 against the predicted 1.920/1.608/1.125/0.800 at $\gst=1.05/1.3/2/3$ (\rep{NEAR_MISS} \S3.4).
This check is for a vacuum near miss at fixed $J=bp^*$, and the formula is post hoc, from
$p_\perp=J/[\ell(1+\sigma\ell/6\mQ)]$ for a rigidly rotating straight string at the turning point.

\subsection{The constant: matched asymptotics for slow ends}\label{app:der:nr}
The constant $K$ in $\rot=\epsp G[\ln(\lmax/\bodd)+K]$ needs the region that the intermediate analysis leaves out. For
slow ends it can be computed in closed form by matching an outer and an inner solution (check \rep{c05}). The units are
$\mu=\sigma=1$, so $\mQ=2$, the relative speed at a crossing is $v_0=2\bst$, $p^*=v_0$, and $\epsilon=q|\Estar_\perp|$
in these units equals $\epsp$.

\paragraph{Setup.} Over one period, measured from an even crossing at $t=0$, the separation and the angular momentum
are
\begin{equation}
  x(t)=t\Bigl(v_0-\frac t2\Bigr),\qquad J(t)=\epsilon\int_0^tx\,dt'=\epsilon\Bigl(\frac{v_0t^2}{2}-\frac{t^3}{6}\Bigr),
  \qquad 0\le t\le2v_0,
\end{equation}
mirrored on $2v_0\le t\le4v_0$. Here $\lmax=v_0^2/2$, $J_1=J(2v_0)=\tfrac23\epsilon v_0^3$ and
$b=J_1/v_0=\tfrac23\epsilon v_0^2$. Choose a matching radius $R_m$ with $b\ll R_m\ll\lmax$. The axis is a line. At the
even crossing the relative vector passes through zero head-on ($J=0$) and the angle $\theta=z/x$ is continuous. At the
odd crossing the relative vector swings around the other end by $\approx\pi$, of which $\pi$ is the trivial reversal.
So the rotation of the axis per period is
\begin{equation}
  \rot=I_{\rm out}(R_m)+\Delta\varphi_{\rm in}(R_m)-\pi .
\end{equation}

\paragraph{Outer region ($|x|>R_m$).} Here $\dot\theta=J/x^2$ (item~3 of Section~\ref{app:der:linear}). The integrand
simplifies to
\begin{equation}
  \frac{J}{x^2}=\frac{\epsilon(v_0/2-t/6)}{(v_0-t/2)^2}=\frac{2\epsilon}{3}\,\frac{v_0+w}{w^2},\qquad w\equiv2v_0-t .
\end{equation}
It is regular at the even crossing ($t\to0$) and has the $1/w^2$ and $1/w$ singularities at the odd one. Stop the
integral at $w=\delta$, where $|x|=R_m$, and double it for the mirror half:
\begin{equation}
  I_{\rm out}=2\cdot\frac{2\epsilon}{3}\int_\delta^{2v_0}\Bigl(\frac{v_0}{w^2}+\frac1w\Bigr)dw
  =\frac{4\epsilon}{3}\Bigl[\frac{v_0}{\delta}-\frac12+\ln\frac{2v_0}{\delta}\Bigr].
\end{equation}
With $R_m=\delta(v_0-\delta/2)$, i.e.\ $v_0/\delta=v_0^2/R_m-\tfrac12+O(R_m)$ and $\ln(2v_0/\delta)=\ln(2v_0^2/R_m)+O(R_m)$, and
with $\tfrac43\epsilon=2b/v_0^2$:
\begin{equation}
  I_{\rm out}(R_m)=\frac{2b}{R_m}+\frac{2b}{v_0^2}\Bigl[\ln\frac{2v_0^2}{R_m}-1\Bigr]+O\!\Bigl(\frac{R_m}{v_0^2}\Bigr).
  \label{eq:D:Iout}
\end{equation}
The $1/w$ term, and hence the logarithm, comes from the deceleration of the relative motion: $x$ is not linear in $w$.

\paragraph{Inner region ($r<R_m$).} Here the field is negligible: it changes $J$ by a fraction $O(R_m^2/v_0^4)$. The
motion is therefore a passage through the central potential $V=\sigma r=r$ at fixed $J=J_1$ and energy $v_0^2/2$, with
$V\ll v_0^2$. The angle swept between entering and leaving radius $R_m$ follows from the radial action, as in the
standard treatment of orbit precession (Landau--Lifshitz, \emph{Mechanics}, \S\S14--15):
\begin{equation}
  \Delta\varphi_{\rm in}=-\frac{\partial S_r}{\partial J},\qquad
  S_r=2\int_{r_{\min}}^{R_m}p_r\,dr,\qquad p_r=\sqrt{v_0^2-2r-J^2/r^2}.
\end{equation}
At zeroth order in $V$, $p_r=\sqrt{v_0^2-J^2/r^2}$ and $\Delta\varphi_0=\pi-2\arcsin(b/R_m)\simeq\pi-2b/R_m$, which is a
straight line at impact parameter $b$. At first order, $p_r\simeq\sqrt{v_0^2-J^2/r^2}-r/\sqrt{v_0^2-J^2/r^2}$, so
\begin{equation}
  S_1=-\frac{2}{v_0}\int_b^{R_m}\frac{r^2\,dr}{\sqrt{r^2-b^2}}
     =-\frac{1}{v_0}\Bigl[R_m\sqrt{R_m^2-b^2}+b^2\ln\frac{R_m+\sqrt{R_m^2-b^2}}{b}\Bigr]
     \simeq-\frac{1}{v_0}\Bigl[R_m^2-\frac{b^2}{2}+b^2\ln\frac{2R_m}{b}\Bigr].
\end{equation}
Since $J=bv_0$, $\Delta\varphi_1=-(1/v_0)\,\partial S_1/\partial b$, and
\begin{equation}
  \Delta\varphi_{\rm in}(R_m)=\pi-\frac{2b}{R_m}+\frac{2b}{v_0^2}\Bigl[\ln\frac{2R_m}{b}-1\Bigr].
  \label{eq:D:inner}
\end{equation}

\paragraph{Matching.} In the sum the $2b/R_m$ terms cancel and $R_m$ drops out, as it must:
\begin{equation}
  \rot_{\rm NR}=\frac{2b}{v_0^2}\Bigl[\ln\frac{4v_0^2}{b}-2\Bigr]
  =\frac43\epsp\Bigl[\ln\frac{6}{\epsp}-2\Bigr]
  =\frac43\epsp\Bigl[\ln\frac{\lmax}{\bodd}+3\ln2-2\Bigr],
  \label{eq:D:rotNR}
\end{equation}
using $4v_0^2/b=6/\epsp$ and $\lmax/b=3/(4\epsp)$. This is Eq.~\eqref{eq:rotNR}, with $K(1)=3\ln2-2=0.0794$. Two
consistency checks:
\begin{itemize}
\item The log coefficient $2b/v_0^2$ equals $\sigma b/(p^*\bst)$ (with $p^*=v_0$, $\bst=v_0/2$). This is the slow-end
  limit of the intermediate-region result~\eqref{eq:D:dprel}, obtained here by a different route.
\item The same calculation with constant $J$ and no field gives the vacuum rotation per crossing,
  $\psi=(b/\lmax)[\ln(\lmax/b)+3\ln2-2]$ (check \rep{c09}). This is measured as $K_{\rm vac}=0.0791\pm0.0003$.
\end{itemize}
The split between the two crossings is $\psi_{\rm even}=\tfrac23+\tfrac43\ln2=1.591$ (in units of $\epsp$; a smooth
torque rotation with no logarithm) and $\psi_{\rm odd}=\tfrac43[\ln(1/\epsp)+\ln3-\tfrac52]$. The two-solver numerics give
1.606 and $1.318\ln(1/\epsp)-1.91$ at $\gst=1.05$ (\rep{NUMERICS.md} M2b; compared in check \rep{c09}).

\paragraph{Relativistic constant.} For relativistic ends the inner problem is the full NG string around a near miss,
which was not solved analytically. $K(\gst)$ is obtained numerically: an exact causal string with the pair at rest in a
pure $\Estar\perp\hat n$ field, $\epsp=10^{-6}$--$10^{-3}$, axis read from the apocentres, and $G$ fixed to
Eq.~\eqref{eq:D:G}. The result is $K=+0.056/-0.049/-0.175/-0.309/-0.569$ at $\gst=1.05/1.28/1.6/2/3$, with a second
solver agreeing at $\gst=1.3$ and 2 (\EST{}).

\subsection{Accumulation and the thickness law}\label{app:der:law}
\paragraph{Coherence.} $\mathbf J_{\rm odd}$ points the same way in every period, because $\hat n$ and $\hat e_\perp$
change only by $O(\Theta)$. Every odd-crossing rotation therefore turns the axis toward $+\Estar_\perp$, and after
$N_{\rm odd}$ odd crossings $\Theta=N_{\rm odd}\rot$. There is no random walk. In the tracker $\Theta\lesssim10^{-2}$,
so the rotations add linearly.

\paragraph{Projection on the hit plane.} In the vacuum the hits lie on the plane through the vertex spanned by
$\hat\beta_{\rm cm}$ and $\hat n$, with normal $\hat N\propto\hat n\times\hat\beta_{\rm cm}$. A boost along
$\hat\beta_{\rm cm}$ leaves components perpendicular to $\hat\beta_{\rm cm}$ unchanged, so the CM and lab planes have
the same $\hat N$. Only the part of the rotation along $\hat N$ moves hits out of the plane. Its share is
$f_N=|\hat e_\perp\cdot\hat N|$. Since $\hat n\cdot\hat N=0$, one has $f_N\sin\alpha=|\hat E^*\cdot\hat N|$: the
out-of-plane rotation is driven by the out-of-plane component of the CM electric field,
\begin{equation}
  f_N\rot=\epsB\,\gamma_{\rm cm}\beta_T\,|\hat E^*\cdot\hat N|\,G(\gst)\Bigl[\ln\frac{\lmax}{\bodd}+K\Bigr].
  \label{eq:D:fNrot}
\end{equation}
For a transverse pair ($\hat\beta_{\rm cm}=\hat x$, $\mathbf B=B\hat z$, $\hat n=(\cos\theta^*,\sin\theta^*\cos\phi^*,
\sin\theta^*\sin\phi^*)$) one has $\hat E^*=-\hat y$ and $\hat N=(0,\sin\phi^*,-\cos\phi^*)$. Hence $\sin\alpha=
\sqrt{1-\sin^2\theta^*\cos^2\phi^*}$ and $|\hat E^*\cdot\hat N|=|\sin\phi^*|$. Pairs whose axis lies in the bending plane
($\phi^*=0$) rotate within their own plane and acquire almost no thickness. This angle drives most of the
event-to-event spread of $\DsB$.

\paragraph{Number of periods.} The pair reaches lab radius $R$ at lab time $R/\beta_T$, i.e.\ CM time
$R/(\gamma_{\rm cm}\beta_T)$. The number of CM periods, and hence of odd crossings, is, with Eq.~\eqref{eq:D:gbT},
\begin{equation}
  N_{\rm odd}(R)\simeq\frac{R}{\gamma_{\rm cm}\beta_T}\cdot\frac{\sigma}{4p^*}=\frac{\sigma R}{2\bst p_T}=\Nosc(R).
  \label{eq:D:Nosc}
\end{equation}

\paragraph{The plane fit sees a twist.} The out-of-plane displacement of a hit is $h=\pm\tfrac{x}{2}\,\Theta\,f_N$. Here
$\pm x/2$ is the in-plane lever of the hit, which oscillates between $-\lmax/2$ and $\lmax/2$, and $\Theta$ grows
linearly with the radius $R_k$ of the hit. So $h=\kappa\,u\,R_k$ with $u=\pm x/2$ and
$\kappa=f_N\rot\Nosc(R_{\rm out})/R_{\rm out}$: a helicoid. The fitted plane through the vertex can tilt about
$\hat\beta_{\rm cm}$ and absorb the mean twist $\kappa\langle R\rangle u$. Since $u$ is symmetric about zero (the two
quirks sit at $\pm u$), it cannot absorb anything else. The residual is $\kappa^2\langle u^2\rangle{\rm Var}(R)$, so
\begin{equation}
  \DsB\simeq g\,f_N\,\rot\,\Nosc(R_{\rm out})\,\frac{\lmax}{2},\qquad
  g\simeq\frac{x_{\rm rms}}{\lmax}\cdot\frac{\sigma_R}{R_{\rm out}}=0.73\times0.32=0.23 ,
  \label{eq:D:g}
\end{equation}
where $x_{\rm rms}/\lmax=0.73$ is the RMS of $x(t)/\lmax$ over a half-period (0.69--0.73 for $\gst=1$--2), and
$\sigma_R/R_{\rm out}=0.32$ is the RMS spread of the ten barrel radii divided by the outer radius $108.3\cm$. The full
rotating-yo-yo model (next paragraph) gives a median $g=0.219$--0.225 (16--84\%: 0.18--0.26) at $\mQ=200$--$1000\GeV$ and
$\Lam=1$, $3\keV$, independent of $\mQ$ and $\Lam$ (check \rep{c16}, added for this appendix; \EST{}).

\paragraph{The law.} Insert Eqs.~\eqref{eq:D:fNrot}, \eqref{eq:D:Nosc}, $\lmax=2(\gst-1)\mQ/\sigma$ and
Eq.~\eqref{eq:D:gbT} into Eq.~\eqref{eq:D:g}. The factors of $\mQ$ and $p_T$ cancel, because $\gamma_{\rm cm}\beta_T\mQ/p_T=1/(2\gst)$:
\begin{equation}
  \DsB\simeq\frac{qBR_{\rm out}}{\sigma}\,
  \underbrace{g\,|\hat E^*\cdot\hat N|\,\frac{G(\gst)(\gst-1)}{4\rho}}_{F}\,
  \Bigl[\ln\frac{1}{\epsB}+C\Bigr],
  \label{eq:D:law}
\end{equation}
which is Eq.~\eqref{eq:law}. The logarithm was rewritten using $\epsp=\epsB\gamma_{\rm cm}\beta_T\sin\alpha$ and
Eq.~\eqref{eq:D:bodd}:
\begin{equation}
  \ln\frac{\lmax}{\bodd}+K=\ln\frac1\epsB
  \underbrace{-\ln(\gamma_{\rm cm}\beta_T\sin\alpha)+\ln\frac{\gst-1}{\gst-\operatorname{asinh}\rho/\rho}+K(\gst)}_{C}.
  \label{eq:D:C}
\end{equation}
$F$ depends only on the CM velocity $\gst$, on angles and on the layer geometry, not on $\mQ$ or $\Lam$ at fixed
velocities. $C$ depends on $\mQ$ only through $\gamma_{\rm cm}\beta_T=p_T/M$, and $B$ and $\Lam$ enter only through
$\epsB$. For a slow pair $F\to g|\hat E^*\cdot\hat N|\bst/6$, so $\DsB$ vanishes linearly with $\bst$: $\lmax\propto{\bst}^2$
and $\Nosc\propto1/\bst$. Equation~\eqref{eq:D:law} is a transparent estimate. The predictions of the note (Table~\ref{tab:c},
and $\DsB^{\rm model}$ and $c_{\rm ev}$ in Table~\ref{tab:p1}) come from the \emph{rotating-yo-yo model}, which evaluates the same physics without the
twist approximation. It takes the exact vacuum yo-yo in the CM frame, steps its axis by $\rot$ toward $\hat e_\perp$ at
each odd crossing, boosts it to the lab, records the hits on the ten barrel layers and fits the plane (check \rep{c06}).
Nothing is fitted to the exact string. The model omits the non-secular first-order offsets (the miss distance itself and
the static polarisation of the string by $\Estar$), the effect of $\Bstar$ (Section~\ref{sec:string:valid}) and the
spreading of the rotation over the half-period at $\gst>2$. Its accuracy against the exact string is 1--4\% for
$\gst\lesssim1.6$.

\paragraph{The slope $c$ and the exponents.} Per event, $F$ and $C$ do not depend on $B$ or $\Lam$, and $\epsB\propto
B\Lam^{-2}$. So $\ln\DsB=\ln\epsB+\ln[\ln(1/\epsB)+C]+{\rm const}$, and differentiating gives exactly
\begin{equation}
  \frac{d\ln\DsB}{d\ln B}=1-c,\qquad \frac{d\ln\DsB}{d\ln\Lam}=-2+2c,\qquad c=\frac{1}{\ln(1/\epsB)+C}.
\end{equation}
The response per unit field follows without expansion:
$\mathcal R(f)=[\ln(1/\epsB)+\ln(1/f)+C]/[\ln(1/\epsB)+C]=1+c\ln(1/f)$. At $500\GeV$ and $1\keV$, $\ln(1/\epsB)=8.40$, and
the population median $c=0.104$ (Table~\ref{tab:c}) corresponds to $C=1.23$. Locally at $\Bzero$ the thickness then
grows like $B^{0.90}$ and falls like $\Lam^{-1.79}$ rather than $B\Lam^{-2}$; averaged over $\Bzero/80\ldots\Bzero$
the $B$ exponent is 0.914 (Table~\ref{tab:numerics}). For charge $q$, $\epsB\to q\epsB$, which gives the $|q|=2$ ratio $2[1-c\ln2]=1.86$ of
Section~\ref{sec:string:valid}.

\subsection{The scale-free property}\label{app:der:scale}
Scale all coordinates and times by $\lambda$ at fixed momenta. Then $d\tau\to\lambda d\tau$, the world-sheet area
$dA\to\lambda^2dA$, and with $A_\mu$ of a uniform field ($\mathbf A=\mathbf B\times\mathbf x/2$), $A\cdot dx\to\lambda^2A\cdot
dx$. Hence $S[\lambda x;\mQ,\sigma,qB]=\lambda\,S[x;\mQ,\lambda\sigma,\lambda qB]$ (check \rep{c15}). A solution for
$(\sigma,qB)$ therefore maps onto one for $(\sigma/\lambda,qB/\lambda)$. The motion per oscillation depends only on
$\epsB=qB/\sigma$ and the velocities, and $\sigma$ (with $\mQ$) sets the length unit $\lzero=\mQ/\sigma$. At a fixed tracker the only other
$\Lam$ dependence is $\Nosc\propto\sigma$ and the lever $\lzero/R$. This is why $(\Lam\sqrt k,B)$ and $(\Lam,B/k)$ are the
same physics per oscillation (Section~\ref{sec:string:valid}).

\subsection{Conditions and status of each step}\label{app:der:status}
The derivation assumes: (i) $\epsp\ll1$ (true by $\gtrsim3.5$ orders of magnitude); (ii) $\Theta\ll1$ over the tracker;
(iii) first order in $\Estar_\parallel$, which only makes successive half-periods unequal; (iv) classical NG at the odd
crossing, $\bodd\gg1/\Lam$, i.e.\ $\Lam\lesssim10$--$26\keV$; (v) $\Bstar$ small in the inner region of the crossing,
$\kB\lesssim4$--6, beyond which its non-linear effect is measured but not derived; (vi) for the hit-level law, the
pair turns around well inside the tracker, median $\lmax/\RIT\lesssim0.6$ (found empirically,
Section~\ref{sec:p1:law}); and (vii) the rotation is localised at the odd crossing, $\gst\lesssim1.6$--2 (rotating-yo-yo
model accurate to 1--4\%; 10--35\% off for $\gst>2$).

\begin{table}[htbp]
\centering
\footnotesize
\caption{Status of each step of the P1 derivation. ``Derived'' means derived analytically and, where a check is listed, checked by sympy and a
limiting case; ``verified'' means confirmed by exact-string numerics (two solvers where stated in
Table~\ref{tab:numerics}); \EST{} means a numerically determined constant or a model step.}
\label{tab:derivation_status}
\begin{tabularx}{\linewidth}{>{\raggedright\arraybackslash}p{0.30\linewidth}>{\raggedright\arraybackslash}Xl}
\toprule
step & status & check\\
\midrule
fields in the CM frame, Eq.~\eqref{eq:D:fields} & exact & ---\\
vacuum yo-yo, half-period integral, Eqs.~\eqref{eq:D:x}--\eqref{eq:D:I} & exact (derived) & c01\\
rigid-rotation lemma & derived & ---\\
torque law, no linear $\Bstar$ effect, Eq.~\eqref{eq:D:torque} & derived & c01 (4)\\
miss distance $\bodd$, Eq.~\eqref{eq:D:bodd} & derived; verified to $10^{-4}$ (c03), $\le10^{-3}$ (two solvers) & c01, c03\\
failure of linear theory only at odd crossings & derived; linear to $10^{-4}$ before the first crossing & c01, c02\\
intermediate NG solution, force, Eqs.~\eqref{eq:D:wedge}--\eqref{eq:D:Fperp} & derived; slope $\bst/\operatorname{artanh}\bst$ verified to $\le0.8\%$ ($\gst\le2$), 3\% at $\gst=3$ & c02, \rep{NUMERICS.md} M2a\\
log coefficient $G$, Eq.~\eqref{eq:D:G} & derived; verified to $\le0.03\%$ (c03), 0.1--0.2\% (two solvers) & c01 (6), c03\\
$K(1)=3\ln2-2$, Eq.~\eqref{eq:D:rotNR} & derived; verified (field: 0.07945; vacuum $K_{\rm vac}=0.0791\pm0.0003$) & c04, c05, c09\\
$K(\gst)$ for relativistic ends & numerical (\EST{}); two solvers at $\gst=1.3$, 2 & c03, c08\\
coherence $\Theta=N_{\rm odd}\rot$, projection $f_N$ & derived (first order in $\Theta$) & c13\\
twist factor $g\simeq0.23$, Eq.~\eqref{eq:D:g} & estimate; model gives 0.22 & c16\\
thickness law, $C$, Eqs.~\eqref{eq:D:law}--\eqref{eq:D:C} & scaling derived; $F$ from the rotating-yo-yo model (\EST{}, 1--4\% for $\gst\lesssim1.6$) & c06, c07\\
$c$, exponents, $\mathcal R(f)$ & exact consequences of the law; verified at $1\keV$ & c07, \rep{NUMERICS.md}\\
scale-free property & exact & c15\\
effect of $\Bstar$ on $\rot$ & numerical only; mechanism not derived & c12, c14\\
\bottomrule
\end{tabularx}
\end{table}
